\documentclass[11pt,a4paper]{article}
% ===============================================================
% Convenio ATLAS II - Provincia de Córdoba - Componente Solar
% Cuarto Informe de Avance: Modelo Satelital Híbrido, Potencial
% Fotovoltaico y Atlas Solar Territorial a 2 km
% ===============================================================

\usepackage[utf8]{inputenc}
\usepackage[spanish,es-nodecimaldot]{babel}
\usepackage[margin=2.5cm]{geometry}
\usepackage{graphicx}
\usepackage{amsmath,amssymb}
\usepackage{booktabs}
\usepackage{array}
\usepackage{multirow}
\usepackage{siunitx}
\usepackage{hyperref}
\usepackage{xcolor}
\usepackage{float}
\usepackage{caption}
\usepackage{subcaption}
\usepackage{fancyhdr}

\hypersetup{
    colorlinks=true,
    linkcolor=blue!70!black,
    citecolor=blue!70!black,
    urlcolor=blue!70!black
}

\pagestyle{fancy}
\fancyhf{}
\rhead{\small Convenio ATLAS II — Componente Solar}
\lhead{\small Cuarto Informe de Avance}
\rfoot{\small Página \thepage}

\begin{document}

% ===============================================================
% PORTADA
% ===============================================================
\begin{titlepage}
\centering

\vspace*{-2cm}

%\includegraphics[height=1.8cm]{figures/logo-fcefyn.png}
\hspace{1.5cm}
%\includegraphics[height=2.2cm]{figures/logo_ohmc.jpeg}

\vspace{4.5cm}

{\fontsize{30}{36}\selectfont \textbf{Convenio ATLAS II} \par}
\vspace{0.4cm}
{\fontsize{26}{32}\selectfont \textbf{Provincia de Córdoba} \par}

\vspace{0.8cm}

{\fontsize{22}{28}\selectfont \textcolor{gray}{\textbf{Componente Solar}} \par}

\vspace{2.5cm}

{\large
Facultad de Ciencias Exactas, Físicas y Naturales (FCEFyN) \par}
\vspace{0.2cm}
{\large
Universidad Nacional de Córdoba (UNC) \par}
\vspace{0.2cm}
{\large
Laboratorio de Hidrometeorología (LHM) / OHMC \par}

\vspace{3.5cm}

{\normalsize \textbf{Cuarto Informe de Avance: Modelo Satelital Híbrido, Potencial Fotovoltaico y Atlas Solar Territorial a 2 km} \par}
\vspace{0.8cm}
{\normalsize Septiembre 2026 \par}

\end{titlepage}

\newpage

% ===============================================================
% RESUMEN EJECUTIVO
% ===============================================================
\section*{Resumen Ejecutivo}



% ===============================================================
% CUARTO INFORME DE AVANCE
% ATLAS SOLAR TERRITORIAL Y POTENCIAL FOTOVOLTAICO DE CÓRDOBA
% ===============================================================

\tableofcontents
\newpage


% ===============================================================
% 1. INTRODUCCIÓN
% ===============================================================

\section{Introducción}

\subsection{Objetivo del Cuarto Informe de Avance}
\subsection{Alcance de esta etapa del proyecto}
\subsection{Síntesis de los desarrollos previos}
\subsection{Productos desarrollados en esta etapa}


% ===============================================================
% 2. CONSTRUCCIÓN METODOLÓGICA DEL ATLAS
% ===============================================================

\section{Construcción Metodológica del Atlas Solar Territorial}

\subsection{Grilla territorial de la Provincia de Córdoba}
\subsection{Resolución espacial y dominio provincial}
\subsection{Regionalización solar del territorio}

\subsection{Asignación de modelos regionales}
\subsubsection{Asignación rígida: Modo Hard}
\subsubsection{Ponderación continua: Modo Blended}

\subsection{Análisis de sensibilidad de la regionalización}
\subsubsection{Comparación Hard vs. Blended}
\subsubsection{Continuidad espacial en las fronteras regionales}
\subsubsection{Impacto sobre las métricas provinciales}

\subsection{Configuración territorial adoptada}
\subsection{Arquitectura de inferencia territorial}
\subsection{Agregación temporal y construcción de climatologías}
\subsection{Control de cobertura y consistencia}
\subsection{Validación numérica de la cadena territorial}


% ===============================================================
% 3. RECURSO SOLAR TERRITORIAL
% ===============================================================

\section{Distribución Territorial del Recurso Solar}

\subsection{Irradiación Global Horizontal (GHI)}
\subsection{Irradiación Directa Normal (DNI)}
\subsection{Irradiación Difusa Horizontal (DHI)}
\subsection{Horas Sol Pico (PSH)}
\subsection{Variabilidad espacial del recurso solar}
\subsection{Variabilidad estacional: verano e invierno}
\subsection{Diferencias regionales dentro de la provincia}


% ===============================================================
% 4. MODELADO Y POTENCIAL FOTOVOLTAICO TERRITORIAL
% ===============================================================

\section{Modelado y Potencial Fotovoltaico Territorial}

\subsection{Cadena física de conversión del recurso solar}

\subsubsection{Descomposición de GHI mediante el modelo de Erbs}
\subsubsection{Transposición al plano mediante el modelo de Perez}
\subsubsection{Modelo térmico de módulo mediante NOCT}
\subsubsection{Conversión DC/AC mediante PVWatts}
\subsubsection{Pérdidas del sistema y supuestos tecnológicos}

\subsection{Configuraciones fotovoltaicas analizadas}
\subsubsection{Sistema fijo a 30° Norte}
\subsubsection{Seguidor horizontal en un eje Norte-Sur}

\subsection{Indicadores energéticos}
\subsubsection{Irradiación en plano del generador}
\subsubsection{Rendimiento específico anual}
\subsubsection{Factor de capacidad}

\subsection{Distribución territorial del potencial fotovoltaico}
\subsubsection{Potencial en configuración fija}
\subsubsection{Potencial con seguimiento en un eje}

\subsection{Identificación de zonas de mayor y menor potencial eléctrico}


% ===============================================================
% 5. GANANCIA POR SEGUIMIENTO Y ANÁLISIS REGIONAL
% ===============================================================

\section{Ganancia por Seguimiento y Análisis Regional}

\subsection{Definición de la ganancia relativa por seguimiento}

\subsection{Distribución espacial de la ganancia}

\subsection{Relación entre recurso solar y ganancia de seguimiento}

\subsection{Influencia de la componente directa de la radiación}

\subsection{Comportamiento estacional de la ganancia}

\subsection{Análisis por regiones solares}

\subsubsection{Región Serrana}
\subsubsection{Región Centro-Sur}
\subsubsection{Región Este}
\subsubsection{Región Norte}

\subsection{Comparación regional}
\subsubsection{Recurso solar}
\subsubsection{Rendimiento fotovoltaico}
\subsubsection{Factor de capacidad}
\subsubsection{Ganancia por seguimiento}

\subsection{Síntesis territorial e implicancias tecnológicas}


% ===============================================================
% 6. PRODUCTOS CARTOGRÁFICOS Y GEOESPACIALES
% ===============================================================

\section{Productos Cartográficos y Base Geoespacial}

\subsection{Atlas cartográfico temático}
\subsection{Rásters Cloud-Optimized GeoTIFF (COG)}
\subsection{GeoPackage provincial}
\subsection{Variables incluidas en los productos geoespaciales}
\subsection{Resolución espacial, sistema de referencia y metadatos}
\subsection{Aplicaciones de los productos generados}


% ===============================================================
% 7. CONCLUSIONES
% ===============================================================

\section{Conclusiones}

\subsection{Distribución provincial del recurso solar}
\subsection{Potencial de generación fotovoltaica}
\subsection{Beneficio territorial del seguimiento en un eje}
\subsection{Principales contrastes regionales}
\subsection{Productos generados para el Atlas Solar}


% ===============================================================
% ANEXOS
% ===============================================================

\appendix

\section{Especificaciones del Sistema Fotovoltaico Modelado}
\section{Formulación Matemática de los Modelos Físicos}
\section{Tablas Estadísticas Provinciales y Regionales}
\section{Especificaciones Técnicas de los Productos Geoespaciales}
\section{Validaciones Numéricas y Pruebas de Consistencia del Pipeline}

\end{document}